% =============================================================================
%  SAMBP Technical Report TR-86/2028
%  IEEE 118-Bus Large-Scale Monte Carlo Protection Study
%  N = 50,000 LHS trials | 5 protection elements | 4 IBR penetration levels
% =============================================================================
\documentclass[11pt,a4paper]{article}

\usepackage[T1]{fontenc}
\usepackage[utf8]{inputenc}
\usepackage[margin=25mm,top=28mm,bottom=28mm]{geometry}
\usepackage{amsmath,amssymb,amsthm}
\usepackage{booktabs,array,multirow,longtable}
\usepackage{graphicx}
\usepackage[table,dvipsnames]{xcolor}
\usepackage[hidelinks,bookmarks=true]{hyperref}
\usepackage{caption,subcaption}
\usepackage{parskip}
\usepackage{siunitx}
\usepackage{pifont}
\usepackage[numbers,sort&compress]{natbib}
\usepackage{fancyhdr}
\usepackage{mdframed}
\usepackage{float}
\usepackage{microtype}

\definecolor{sambpblue}{RGB}{0,70,140}
\definecolor{okgreen}{RGB}{0,130,60}
\definecolor{alertred}{RGB}{180,0,0}
\definecolor{warnoran}{RGB}{200,100,0}
\definecolor{lightgrey}{RGB}{240,240,240}
\definecolor{passgreen}{RGB}{230,255,230}
\definecolor{failred}{RGB}{255,230,230}

\newcommand{\pu}{\,\text{pu}}
\newcommand{\ms}{\,\text{ms}}
\newcommand{\cmark}{\textcolor{okgreen}{\ding{51}}}
\newcommand{\xmark}{\textcolor{alertred}{\ding{55}}}

\mdfdefinestyle{finding}{
  backgroundcolor=passgreen, linecolor=okgreen, linewidth=1pt,
  innertopmargin=6pt, innerbottommargin=6pt,
  frametitle={\textbf{Finding}}, frametitlebackgroundcolor=okgreen!20}
\mdfdefinestyle{alert}{
  backgroundcolor=failred, linecolor=alertred, linewidth=1pt,
  innertopmargin=6pt, innerbottommargin=6pt,
  frametitle={\textbf{Key Result}}, frametitlebackgroundcolor=alertred!20}

\pagestyle{fancy}
\fancyhf{}
\lhead{\textcolor{sambpblue}{\textbf{SAMBP TR-86/2028}}}
\rhead{\textcolor{gray}{\small IEEE 118-Bus Monte Carlo}}
\cfoot{\thepage}

\title{%
  \textbf{\textsf{\color{sambpblue}SAMBP Technical Report TR-86/2028}}\\[6pt]
  \Large IEEE 118-Bus Large-Scale Monte Carlo\\
  Protection Validation Study\\[4pt]
  \normalsize Latin Hypercube Sampling ($N = 50{,}000$) over 6 Dimensions\\
  5 Protection Elements: 87L / 21 / 46 / 51 / 67}
\author{SAMBP Research Group\\[2pt]
  \small Smart Adaptive Multi-Bus Protection Research, IIT Madras}
\date{April 2028\quad\textbar\quad Chapters 4 \& 5 (87L + Distance)\\
  Prerequisite: TR-70 (Meshed IBR), TR-85 (CLPU)}

\begin{document}
\maketitle
\tableofcontents
\clearpage

% =============================================================================
\section{Background and Objectives}
\label{sec:background}
% =============================================================================

\subsection{Motivation}

All prior SAMBP technical reports (TR-01 through TR-85) validated protection
algorithms on single-feeder or small multi-bus networks (up to 39 buses).
A critical open question is whether the same algorithms scale to
\emph{transmission-class meshed networks} with hundreds of buses and
heterogeneous IBR penetration.

The IEEE 118-bus test system \cite{Christie1993} is the standard large-scale
benchmark in power systems protection research.  It comprises 118 buses,
186 transmission lines, 54 generators, and 91 load points, offering
representative impedance diversity, topological complexity, and multiple
voltage levels.

\subsection{Hypotheses}

Two quantitative hypotheses are tested:

\begin{mdframed}[style=finding]
\textbf{H1 (Primary):} The SAMBP 87L line differential relay achieves an
agree-rate $\geq 97\%$ across all bus/fault/load combinations when IBR
penetration is $\leq 50\%$.
\end{mdframed}

\begin{mdframed}[style=finding]
\textbf{H2 (Secondary):} All five SAMBP protection elements (87L, 21, 46,
51, 67) achieve agree-rate $\geq 90\%$ at full IBR penetration (100\%).
\end{mdframed}

\subsection{Scope}

\begin{itemize}
  \item \textbf{Network:} IEEE 118-bus, pandapower \texttt{case118()}
  \item \textbf{Trials:} $N = 50{,}000$ Latin Hypercube Sampling (LHS)
  \item \textbf{LHS dimensions:} 6 (bus index, fault type, fault resistance,
        IBR penetration, IBR mode, load level)
  \item \textbf{Protection elements:} 87L, 21, 46, 51, 67
  \item \textbf{IBR penetration:} 0\%, 30\%, 50\%, 100\%
  \item \textbf{IBR modes:} GFL, GFM, mixed
  \item \textbf{Fault types:} 3PH, SLG, DLG, LL
  \item \textbf{Fault resistance:} Log-normal, $R_f \in [0, 20]\,\Omega$
  \item \textbf{Load level:} Uniform, $\lambda_L \in [0.7, 1.3]$
\end{itemize}

% =============================================================================
\section{Methodology}
\label{sec:method}
% =============================================================================

\subsection{IEEE 118-Bus IBR Overlay}

The standard \texttt{case118()} network is augmented by replacing the
top-$N$ generators (sorted by rated MW in descending order) with virtual
IBR sources until the target penetration fraction is reached:

\begin{equation}
  \alpha_{\text{IBR}} = \frac{\sum_{g \in \mathcal{G}_{\text{IBR}}} P_g^{\max}}
  {\sum_{g \in \mathcal{G}} P_g^{\max}},
  \label{eq:ibr_pen}
\end{equation}

where $\mathcal{G}_{\text{IBR}} \subseteq \mathcal{G}$ is the subset of
generators replaced.  GFL inverters are modelled with current limit
$I_{\max}^{\text{GFL}} = 1.2\pu$; GFM inverters with
$I_{\max}^{\text{GFM}} = 1.5\pu$.

\subsection{Analytical Fault Model}

For computational tractability at $N = 50{,}000$, fault currents are computed
analytically from the Thevenin equivalent at each bus rather than via
full power-flow iteration.  The Thevenin impedance at bus $k$ is estimated
from the minimum connected line impedance:

\begin{equation}
  Z_{\text{th},k} = \min_{(i,k) \in \mathcal{E}} Z_{ik},
  \label{eq:z_th}
\end{equation}

which provides a conservative (worst-case) fault level estimate.

The fault current positive-sequence component for each fault type is:

\begin{align}
  \text{3PH:}\quad & I_{\text{pos}} = \frac{V_{\text{pre}}}{|Z_{\text{th}} + Z_f|} \\
  \text{SLG:}\quad & I_{\text{pos}} = 0.577\,\frac{V_{\text{pre}}}{|Z_{\text{th}} + Z_f|} \\
  \text{DLG:}\quad & I_{\text{pos}} = 0.816\,\frac{V_{\text{pre}}}{|Z_{\text{th}} + Z_f|} \\
  \text{LL:} \quad & I_{\text{pos}} = 0.707\,\frac{V_{\text{pre}}}{|Z_{\text{th}} + Z_f|}
  \label{eq:fault_currents}
\end{align}

\subsection{IBR Current Limiting and Proportional Scaling}

A key finding of this study (detailed in Section~\ref{sec:finding_ineg}) is
that IBR current limiting must be applied \emph{proportionally} to both
positive- and negative-sequence currents.  The corrected model is:

\begin{equation}
  I_{\text{cap}} = I_{\max} \bigl(1 + 0.5(1 - k_{\text{IBR}})\bigr),
  \label{eq:i_cap}
\end{equation}

\begin{equation}
  \text{If } I_{\text{pos}} > I_{\text{cap}}: \quad
  s = \frac{I_{\text{cap}}}{I_{\text{pos}}}, \quad
  I_{\text{pos}} \leftarrow I_{\text{cap}}, \quad
  I_{\text{neg}} \leftarrow s \cdot I_{\text{neg}},
  \label{eq:prop_cap}
\end{equation}

where $s \in (0,1]$ is the proportional scale factor.  This preserves the
symmetrical-component ratio $I_{\text{neg}}/I_{\text{pos}}$ under IBR
limiting, which is required for correct differential relay operation.

\subsection{Latin Hypercube Sampling Design}

LHS partitions each dimension into $N$ equal-probability intervals and
draws one sample per interval per dimension, then randomly combines them.
For $N = 50{,}000$ and 6 dimensions, the LHS design provides excellent
coverage of the parameter space with far fewer samples than a full-factorial
grid ($118 \times 4 \times 4 \times 4 \times 3 \approx 22{,}500$ grid points
would cover only the discrete dimensions).

The continuous dimensions (fault resistance, load level) use the mapping:

\begin{align}
  R_f &= \exp\bigl(\mu_{R_f} + \sigma_{R_f} \cdot \Phi^{-1}(u_3)\bigr)
         \;\text{clipped to } [0,20]\,\Omega, \quad
         \mu_{R_f}=0.5,\; \sigma_{R_f}=1.0 \label{eq:Rf_lognormal}\\
  \lambda_L &= 0.7 + 0.6\, u_6, \quad u_6 \sim U[0,1] \label{eq:load_lhs}
\end{align}

\subsection{Protection Stack}

Five SAMBP protection elements are evaluated for each fault scenario:

\begin{table}[H]
\centering
\caption{SAMBP protection elements evaluated in TR-86.}
\label{tab:elements}
\begin{tabular}{lll}
\toprule
Element & Function & Key parameter \\
\midrule
87L & Line differential   & $I_{\text{op}} - 0.25\,I_{\text{neg}} > 0.10 + 0.03\,f_{\text{GFM}}$ \\
21  & Distance/impedance  & $Z_{\text{app,corr}} < 0.85\,Z_{\text{line}}$ (Zone~1) \\
46  & Neg-seq overcurrent & $I_{\text{neg}} > 0.10\pu$, $I_{\text{neg}}/I_{\text{pos}} > 0.20$ \\
51  & Time-overcurrent    & IEC Standard Inverse, $I_{\text{pu}} = 1.25$ \\
67  & Directional OC      & $I_{\text{pos}} > 0.80\pu$ and $V_{\text{pos}} < 0.85\pu$ \\
\bottomrule
\end{tabular}
\end{table}

The 87L operating criterion uses the \emph{percentage-bias} formulation
derived in TR-86:

\begin{equation}
  I_{\text{op}} = I_{\text{pos}}, \quad
  I_{\text{res}} = K_{\text{bias}} \cdot I_{\text{neg}}, \quad
  K_{\text{bias}} = 0.25,
  \label{eq:87L_criterion}
\end{equation}

\begin{equation}
  \text{TRIP iff} \quad I_{\text{op}} - I_{\text{res}} >
  \underbrace{0.10 + 0.03\,f_{\text{GFM}}}_{\text{pickup threshold}}.
  \label{eq:87L_trip}
\end{equation}

% =============================================================================
\section{Key Finding: Proportional I\texorpdfstring{$_{\text{neg}}$}{neg}
         Capping}
\label{sec:finding_ineg}
% =============================================================================

\begin{mdframed}[style=alert]
\textbf{TR-86 Finding F1 — Uncapped $I_{\text{neg}}$ causes false
87L RESTRAINs at IBR penetration $\geq 30\%$:}

The naive IBR model caps $I_{\text{pos}}$ at $I_{\text{cap}}$ but leaves
$I_{\text{neg}}$ at its uncapped (pre-limiting) value.  For high-impedance
buses where $I_{\text{fault}} \gg I_{\text{cap}}$, this results in
$I_{\text{neg}} \gg I_{\text{pos}}$, making the bias term
$K_{\text{bias}} \cdot I_{\text{neg}}$ larger than $I_{\text{op}}$ and
causing the differential relay to restrain on a genuine internal fault.
\end{mdframed}

\subsection{Pre-Fix Diagnosis}

Before applying Eq.~(\ref{eq:prop_cap}), the 87L agree-rate at 30\% IBR
penetration was only 59.2\% (DLG faults: 38.4\%) despite agree-rates of
100\% at 0\% IBR.  Diagnosis showed that for buses with small $Z_{\text{th}}$
(low-impedance buses), $I_{\text{fault}} \approx 12$--$15\pu$, which after
capping gives $I_{\text{pos}} \approx 1.62\pu$ but $I_{\text{neg,raw}}
\approx 7.1\pu$ (DLG ratio 0.55).  The bias term $0.25 \times 7.1 = 1.78\pu$
exceeds $I_{\text{op}} = 1.62\pu$, causing a spurious RESTRAIN.

\subsection{Post-Fix Results}

After applying the proportional scaling of Eq.~(\ref{eq:prop_cap}),
$I_{\text{neg}}$ is scaled to $0.55 \times I_{\text{cap}} = 0.89\pu$ for
DLG faults.  The bias term becomes $0.25 \times 0.89 = 0.22\pu$, well
below $I_{\text{op}} = 1.62\pu$, and the differential relay correctly
TRIPs.  The 87L agree-rate at 30\% IBR recovers to \textbf{100.00\%}.

% =============================================================================
\section{Results}
\label{sec:results}
% =============================================================================

\subsection{Overall Agree-Rate per Protection Element}

\begin{figure}[H]
  \centering
  \includegraphics[width=0.85\textwidth]{/root/phd_thesis/04_code/sambp/digital_twin/reports/figures/TR86/fig1_agree_per_function.pdf}
  \caption{SAMBP TRIP agree-rate per protection element, IEEE 118-bus,
           $N = 50{,}000$ LHS trials.  87L and 67 achieve 100\%; 21, 46,
           51 are degraded at 100\% IBR penetration, pulling the mean below
           90\%.  Dashed red = H1 threshold (97\%); dotted orange = H2
           threshold (90\%).
           \textit{Source: TR-86, mc\_runner\_tr86.py.}}
  \label{fig:agree_per_function}
\end{figure}

Table~\ref{tab:overall} summarises the overall agree-rates across all
$N = 50{,}000$ scenarios (all IBR penetration levels combined).

\begin{table}[H]
\centering
\caption{Overall TRIP agree-rate per protection element.
         $N = 50{,}000$ LHS scenarios, IEEE 118-bus.}
\label{tab:overall}
\begin{tabular}{lccc}
\toprule
Element & Agree-rate & H1 ($\geq 97\%$ at pen $\leq 50\%$) & H2 ($\geq 90\%$ at 100\%) \\
\midrule
87L & 100.00\% & \cmark\ PASS & \cmark\ PASS \\
21  & 25.02\%  & \cmark\ PASS (pen $\leq 50\%$) & \xmark\ FAIL \\
46  & 44.93\%  & \cmark\ PASS (pen $\leq 50\%$) & \xmark\ FAIL \\
51  & 39.97\%  & \cmark\ PASS (pen $\leq 50\%$) & \xmark\ FAIL \\
67  & 100.00\% & \cmark\ PASS & \cmark\ PASS \\
\bottomrule
\end{tabular}
\end{table}

\subsection{Agree-Rate vs IBR Penetration}

\begin{figure}[H]
  \centering
  \includegraphics[width=0.85\textwidth]{/root/phd_thesis/04_code/sambp/digital_twin/reports/figures/TR86/fig2_agree_vs_penetration.pdf}
  \caption{Agree-rate as a function of IBR penetration for all five
           protection elements.  87L (bold blue) maintains 100\% across
           all penetration levels.  Elements 21, 46, 51 degrade sharply
           beyond 50\% IBR.  The H1 region (green shaded, pen $\leq 50\%$)
           shows 87L safely above the 97\% threshold.
           \textit{Source: TR-86.}}
  \label{fig:agree_vs_pen}
\end{figure}

The complete per-penetration per-element agree-rate matrix is given in
Table~\ref{tab:pen_matrix}.

\begin{table}[H]
\centering
\caption{Agree-rate [\%] by IBR penetration and protection element.
         Cells $\geq 97\%$ are in green; cells $< 90\%$ are in red.}
\label{tab:pen_matrix}
\begin{tabular}{lrrrrr}
\toprule
IBR pen. & 87L & 21 & 46 & 51 & 67 \\
\midrule
\rowcolor{passgreen}
0\%   & 100.0 & 100.0 & 100.0 & 100.0 & 100.0 \\
\rowcolor{passgreen}
30\%  & 100.0 & 100.0 & 100.0 & 100.0 & 100.0 \\
\rowcolor{passgreen}
50\%  & 100.0 & 100.0 & 100.0 & 100.0 & 100.0 \\
\rowcolor{failred}
100\% & 100.0 & 25.0  & 44.9  & 40.0  & 100.0 \\
\bottomrule
\end{tabular}
\end{table}

\subsection{Agree-Rate by Fault Type}

\begin{figure}[H]
  \centering
  \includegraphics[width=0.90\textwidth]{/root/phd_thesis/04_code/sambp/digital_twin/reports/figures/TR86/fig3_agree_by_fault_type.pdf}
  \caption{Grouped bar chart of agree-rate by fault type and protection
           element.  87L achieves 100\% for all four fault types.
           Distance relay (21) degrades most severely for LL faults.
           \textit{Source: TR-86.}}
  \label{fig:agree_by_fault}
\end{figure}

\subsection{Agree-Rate Heatmap}

\begin{figure}[H]
  \centering
  \includegraphics[width=0.70\textwidth]{/root/phd_thesis/04_code/sambp/digital_twin/reports/figures/TR86/fig4_heatmap_21.pdf}
  \caption{Heatmap of agree-rate [\%] for element $\times$ IBR penetration.
           Green = $\geq 97\%$, red = $< 90\%$.  The 100\% IBR column
           reveals a clear gap in distance and overcurrent elements, but
           87L and 67 remain immune.
           \textit{Source: TR-86.}}
  \label{fig:heatmap}
\end{figure}

\subsection{Operate-Time Distributions}

\begin{figure}[H]
  \centering
  \includegraphics[width=0.85\textwidth]{/root/phd_thesis/04_code/sambp/digital_twin/reports/figures/TR86/fig5_operate_times.pdf}
  \caption{Empirical CDF of operate time for TRIP decisions across all
           $N = 50{,}000$ scenarios.  87L (blue) operates entirely within
           ${\leq}\,30\,\mathrm{ms}$.  Distance relay (21) shows a bimodal
           distribution at $\sim25\,\mathrm{ms}$ (Zone~1) and
           $\sim400\,\mathrm{ms}$ (Zone~2).
           \textit{Source: TR-86.}}
  \label{fig:op_times}
\end{figure}

% =============================================================================
\section{Hypothesis Test Summary}
\label{sec:hypotheses}
% =============================================================================

\begin{table}[H]
\centering
\caption{Hypothesis test results.}
\label{tab:hypotheses}
\begin{tabular}{lllc}
\toprule
Hyp. & Statement & Result & Verdict \\
\midrule
H1 & 87L agree $\geq 97\%$ for pen $\leq 50\%$ &
    100\% at 0\%, 30\%, 50\% IBR & \cmark\ PASS \\
H2 & All elements $\geq 90\%$ at 100\% IBR &
    21: 25\%, 46: 45\%, 51: 40\% & \xmark\ FAIL \\
\bottomrule
\end{tabular}
\end{table}

H1 is \textbf{confirmed}: the SAMBP 87L relay satisfies the primary
hypothesis on the full IEEE 118-bus network with $N = 50{,}000$ LHS trials.

H2 \textbf{fails} for elements 21, 46, and 51 at 100\% IBR penetration.
This is a new finding not observable on smaller test networks and motivates
three follow-on TRs:
\begin{itemize}
  \item \textbf{TR-87}: IBR-aware distance relay ($Z_{\text{app}}$
        correction for 100\% converter-dominated networks)
  \item \textbf{TR-88}: Negative-sequence OC (46) redesign for GFM-dominant
        grids
  \item \textbf{TR-89}: Graph Neural Network (GNN) protection coordinator
        trained on 118-bus scenarios
\end{itemize}

% =============================================================================
\section{Validation Statistics}
\label{sec:stats}
% =============================================================================

\begin{table}[H]
\centering
\caption{Monte Carlo run statistics.}
\label{tab:mc_stats}
\begin{tabular}{lr}
\toprule
Parameter & Value \\
\midrule
Total LHS scenarios         & 50{,}000 \\
Total protection decisions  & 250{,}000 \\
Valid decisions             & 250{,}000 (0 errors) \\
Wall time                   & 322\,s (5.4\,min) \\
Throughput                  & $\approx 155$\,scenarios/s \\
Worker processes            & 4 \\
Random seed                 & 42 \\
\midrule
Bus index range             & 0\,--\,117 \\
Fault type distribution     & LHS-uniform across \{3PH, SLG, DLG, LL\} \\
$R_f$ distribution          & Log-normal, $\mu=0.5$, $\sigma=1.0$, clip $[0,20]\,\Omega$ \\
Load level distribution     & Uniform $[0.7, 1.3]$ \\
IBR penetration levels      & \{0\%, 30\%, 50\%, 100\%\} (LHS-stratified) \\
IBR modes                   & \{GFL, GFM, mixed\} (LHS-stratified) \\
\bottomrule
\end{tabular}
\end{table}

\subsection{Agree-Rate Confidence Intervals}

For the 87L element at 30\% IBR penetration, $n = 12{,}500$ trials with
$k = 12{,}500$ TRIPs gives agree-rate $\hat{p} = 1.000$.  The 95\%
one-sided Wilson lower bound is:

\begin{equation}
  p_L = \frac{\hat{p} + \frac{z^2}{2n} - z\sqrt{\frac{\hat{p}(1-\hat{p})}{n}
        + \frac{z^2}{4n^2}}}{1 + \frac{z^2}{n}},
  \quad z = 1.645,
  \label{eq:wilson}
\end{equation}

giving $p_L = 99.97\%$, well above the 97\% H1 threshold.

% =============================================================================
\section{Conclusions and Contributions}
\label{sec:conclusions}
% =============================================================================

\begin{enumerate}
  \item \textbf{H1 confirmed:} SAMBP 87L achieves 100\% agree-rate on the
        IEEE 118-bus network for IBR penetration $\leq 50\%$, validated
        over $N = 50{,}000$ LHS trials spanning all bus locations, fault
        types, fault resistances, and load levels.

  \item \textbf{Proportional $I_{\text{neg}}$ capping (TR-86-F1):} A new
        modelling requirement is identified: IBR current limiters must cap
        $I_{\text{neg}}$ proportionally with $I_{\text{pos}}$ to maintain
        the correct symmetrical-component ratio.  Failure to do so reduces
        87L agree-rate from 100\% to 59\% at 30\% IBR penetration.

  \item \textbf{100\% IBR gap:} Elements 21, 46, and 51 fail H2 at 100\%
        IBR penetration, motivating dedicated redesign in TR-87--TR-89.

  \item \textbf{Speed:} The 118-bus MC framework processes 155
        scenarios/second on 4 CPU cores, enabling rapid what-if studies
        during relay setting optimisation.
\end{enumerate}

\subsection{Agree-Rate Formula — 87L Calibrated Threshold}

The final calibrated 87L criterion for IBR networks (applicable for
penetration $\leq 50\%$) is:

\begin{equation}
  \boxed{
    \text{TRIP iff} \quad
    I_{\text{pos}} - 0.25\,I_{\text{neg}}^{(\text{cap})} \;>\; 0.10 + 0.03\,f_{\text{GFM}}
  }
  \label{eq:87L_final}
\end{equation}

where $I_{\text{neg}}^{(\text{cap})}$ is the proportionally-capped
negative-sequence current from Eq.~(\ref{eq:prop_cap}).

\bigskip
\noindent\textbf{Source code:}
\begin{itemize}
  \item Network model: \texttt{04\_code/sambp/digital\_twin/models/ieee118\_network.py}
  \item LHS generator: \texttt{04\_code/sambp/digital\_twin/models/mc\_scenario\_generator.py}
  \item Protection stack: \texttt{04\_code/sambp/digital\_twin/models/sambp\_protection\_stack.py}
  \item MC runner: \texttt{04\_code/sambp/digital\_twin/models/mc\_runner\_tr86.py}
  \item Raw results: \texttt{04\_code/sambp/digital\_twin/results/tr86\_mc\_results.npz}
\end{itemize}

\bibliographystyle{IEEEtranN}
\bibliography{references86}

\end{document}
